
\section{Accelerated Gradient Descent via Plane Search}

According to legend, the first proof for accelerated gradient descent
is due to Nemirovski in the 70s. The first proof involves a $2$-dimensional
plane search subroutine. Later on this was improved to line search
and finally Nesterov showed how to get rid of line search in 1982.
We change the proof slightly to get rid of all uses of parameters.

\begin{algorithm2e}[H]

\caption{$\mathtt{NemAGD}$}

\SetAlgoLined

\textbf{Input:} Initial point $x^{(1)}\in\Rn$.

\For{$k=1,2,\cdots$}{

$x^{(k)+}\leftarrow\text{argmin}_{x=x^{(k)}+t\nabla f(x^{(k)})}f(x).$

$P^{(k)}\leftarrow x^{(1)}+\mathrm{span}(x^{(k)+}-x^{(1)},\sum_{s=1}^{k}\frac{\nabla f(x^{(s)})}{\|\nabla f(x^{(s)})\|})$.

\tcp{Alternatively, one can use $P^{(k)}\leftarrow x^{(k)}+\mathrm{span}(x^{(k)}-x^{(1)},\nabla f(x^{(k)}),\sum_{s=1}^{k}\frac{\nabla f(x^{(s)})}{\|\nabla f(x^{(s)})\|})$
without defining $x^{(k)+}$.}

$x^{(k+1)}=\arg\min_{x\in P^{(k)}}f(x)$.

}

\end{algorithm2e}
\begin{thm}
Assume that $f$ is convex with $\nabla^{2}f(x)\preceq L\cdot I$
for all $x$. Then, we have that 
\[
f(x^{(k)})-f(x^{*})\lesssim\frac{L\|x^{*}-x^{(1)}\|^{2}}{k^{2}}.
\]
\end{thm}

\begin{proof}
Let $\delta_{k}=f(x^{(k)})-f(x^{*})$. By the convexity of $f$, we
have

\begin{align*}
\delta_{k} & \leq\nabla f(x^{(k)})^{\top}(x^{(k)}-x^{*})\\
 & =\nabla f(x^{(k)})^{\top}(x^{(k)}-x^{(1)})+\nabla f(x^{(k)})^{\top}(x^{(1)}-x^{*}).
\end{align*}
Since $x^{(k)}$ is the minimizer on $P^{(k-1)}$ which contains $x^{(1)}$,
we have $\nabla f(x^{(k)})^{\top}(x^{(k)}-x^{(1)})=0$ and hence
\[
\delta_{k}\leq\nabla f(x^{(k)})^{\top}(x^{(1)}-x^{*}).
\]

Let $\lambda_{t}=\frac{1}{\|\nabla f(x^{(t)})\|}$. Note that
\[
\sum_{k=1}^{T}\lambda_{k}\delta_{k}\leq\left\langle \sum_{k=1}^{T}\frac{\nabla f(x^{(k)})}{\|\nabla f(x^{(k)})\|},x^{(1)}-x^{*}\right\rangle \leq\|x^{*}-x^{(1)}\|_{2}\cdot\norm{\sum_{k=1}^{T}\frac{\nabla f(x^{(k)})}{\|\nabla f(x^{(k)})\|}}_{2}.
\]
Finally, we note that $\nabla f(x^{(k+1)})\perp P^{(k)}$ and hence
$\nabla f(x^{(k+1)})\perp\sum_{s=1}^{k}\frac{\nabla f(x^{(s)})}{\|\nabla f(x^{(s)})\|}$.
Therefore, we have 
\[
\norm{\sum_{k=1}^{T}\frac{\nabla f(x^{(k)})}{\|\nabla f(x^{(k)})\|}}_{2}^{2}=\sum_{k=1}^{T}\norm{\frac{\nabla f(x^{(k)})}{\|\nabla f(x^{(k)})\|}}_{2}^{2}=T.
\]
Hence, we have
\[
\sum_{k=1}^{T}\lambda_{k}\delta_{k}\leq\sqrt{T}\cdot\|x^{*}-x^{(1)}\|_{2}.
\]
Since $x^{(k)+}\in P$, we have that 
\[
\delta_{k+1}=f(x^{(k+1)})-f(x^{*})\leq f(x^{(k)+})-f(x^{*})\leq\delta_{k}-\frac{1}{2L}\|\nabla f(x^{(k)})\|_{2}^{2}.
\]

Hence, we have $\|\nabla f(x^{(k)})\|_{2}^{2}\leq2L(\delta_{k}-\delta_{k+1})$.
This gives
\[
\sum_{k=1}^{T}\frac{\delta_{k}}{\sqrt{\delta_{k}-\delta_{k+1}}}\leq\sqrt{2LT}\cdot\|x^{*}-x^{(1)}\|_{2}
\]
for all $T$. Solving this recursion gives the result.
\end{proof}
\begin{xca}
Solve the recursion omitted at the end of the proof.
\end{xca}

Note that the proof above used only the fact that $\{x^{(1)},x^{(k)+},\sum_{s=1}^{k}\frac{\nabla f(x^{(s)})}{\|\nabla f(x^{(s)})\|}\}\subset P$.
Therefore, one can put extra vectors in $P$ to obtain extra features.
For example, if we use the subspace
\[
P=x^{(k)}+\mathrm{span}(x^{(k)}-x^{(1)},\nabla f(x^{(k)}),\sum_{s=1}^{k}\frac{\nabla f(x^{(s)})}{\|\nabla f(x^{(s)})\|},x^{(k)}-x^{(k-1)}),
\]
then one can prove that this algorithm is equivalently to conjugate
gradient when $f$ is a quadratic function.

\section{Accelerated Gradient Descent\label{sec:Accelerated-Gradient-Descent}}


\subsection{Gradient Mapping}

To give a general version of acceleration, we consider problem of
the form
\[
\min_{x}\phi(x)\defeq f(x)+h(x)
\]
where $f(x)$ is strongly convex and smooth and $h(x)$ is convex.
We assume that we access the function $f$ and $h$ differently via:
\begin{enumerate}
\item Let $\mathcal{T}_{f}$ be the cost of computing $\nabla f(x)$.
\item Let $\mathcal{T}_{h,\lambda}$ be the cost of minimizing $h(x)+\frac{\lambda}{2}\norm{x-c}_{2}^{2}$
exactly.
\end{enumerate}
The idea is to move whatever we can optimize in $\phi$ to $h$ and
hopefully this makes the remaining part of $\phi$, $f$, as smooth
and strongly convex as possible. To make the statement general, we
only assume $h$ is convex and hence $h$ may not be differentiable.
To handle this issue, we need to define an approximate derivative
of $h$ that we can compute.
\begin{defn}
We define the gradient step
\[
p_{x}=\argmin_{y}f(x)+\nabla f(x)^{\top}(y-x)+\frac{L}{2}\norm{y-x}^{2}+h(y)
\]
and the gradient mapping
\[
g_{x}=L(x-p_{x}).
\]

Note that if $h=0$, then $p_{x}=x-\frac{1}{L}\nabla f(x)$ and $g_{x}=\nabla f(x)$.
In general, if $\phi\in\mathcal{C}^{2}$, then we have that
\[
p_{x}=x-\frac{1}{L}\nabla\phi(x)+O(\frac{1}{L^{2}}).
\]
Therefore, we have that $g_{x}=\nabla\phi(x)+O(\frac{1}{L})$. Hence,
the gradient mapping is an approximation of the gradient of $\phi$
that is computable in time $\mathcal{T}_{g}+\mathcal{T}_{h,L}$.

The key lemma we use here is that $\phi$ satisfies a lower bound
defining using $g_{x}$. Ideally, we would love to get a lower bound
as follows:
\[
\phi(z)\geq\phi(x)+g_{x}^{\top}(z-x)+\frac{\mu}{2}\norm{z-x}_{2}^{2}.
\]
But it is WRONG. If that was true for all $z$, then we would have
$g_{x}=\nabla\phi(x)$. However, if $\phi\in\mathcal{C}^{2}$ is $\mu$
strongly convex, then we have
\begin{align}
\phi(z) & \geq\phi(x)+\nabla\phi(x)^{\top}(z-x)+\frac{\mu}{2}\norm{z-x}_{2}^{2}\nonumber \\
 & \geq\phi(x-\frac{1}{L}\nabla\phi(x))+\frac{1}{2L}\norm{\nabla\phi(x)}^{2}+\nabla\phi(x)^{\top}(z-x)+\frac{\mu}{2}\norm{z-x}_{2}^{2}.\label{eq:gradient_lower_bound}
\end{align}
It turns out that this is true and is exactly what we need for proving
gradient descent, mirror descent and accelerated gradient descent.
\end{defn}

\begin{thm}
\label{thm:gradient_mapping_lower_bound}Given $\phi=f+h$. Suppose
that $f$ is $\mu$ strongly convex with $L$-Lipschitz gradient.
Then, for any $z$, we have that
\[
\phi(z)\geq\phi(p_{x})+g_{x}^{\top}(z-x)+\frac{1}{2L}\norm{g_{x}}_{2}^{2}+\frac{\mu}{2}\norm{z-x}_{2}^{2}.
\]
\end{thm}

\begin{proof}
Let $\overline{f}(y)=f(x)+\nabla f(x)^{\top}(y-x)+\frac{L}{2}\norm{y-x}_{2}^{2}$
and $p_{t}=p_{x}+t(z-p_{x})$. Using that $p_{0}$ is the minimizer
of $\overline{f}+h$, we have that 
\[
\overline{f}(p_{0})+h(p_{0})\leq\overline{f}(p_{t})+h(p_{t})\leq\overline{f}(p_{0})+\nabla\overline{f}(p_{0})^{\top}(p_{t}-p_{0})+\frac{L}{2}\norm{p_{t}-p_{0}}^{2}+h(p_{t}).
\]
Hence, we have that
\begin{align*}
0 & \leq\nabla\overline{f}(p_{0})^{\top}(p_{t}-p_{0})+\frac{L}{2}\norm{p_{t}-p_{0}}^{2}+h(p_{t})-h(p_{0})\\
 & \leq t\cdot\left(\nabla\overline{f}(p_{0})^{\top}(z-p_{0})+h(z)-h(p_{0})\right)+\frac{Lt^{2}}{2}\norm{z-p_{0}}^{2}.
\end{align*}
Taking $t\rightarrow0^{+}$, we have 
\[
\nabla\overline{f}(p_{x})^{\top}(z-p_{x})+h(z)-h(p_{x})\geq0
\]

Expanding the term $\nabla\overline{f}(p_{x})$, we have
\[
\nabla f(x)^{\top}(z-p_{x})+L(p_{x}-x)^{\top}(z-p_{x})+h(z)-h(p_{x})\geq0.
\]
Equivalently,
\begin{align*}
h(z) & \geq h(p_{x})+\nabla f(x)^{\top}(p_{x}-z)+L(p_{x}-x)^{\top}(p_{x}-z)\\
 & =h(p_{x})+\nabla f(x)^{\top}(p_{x}-z)+\frac{1}{L}\|g_{x}\|^{2}+L(p_{x}-x)^{\top}(x-z)\\
 & =h(p_{x})+\nabla f(x)^{\top}(p_{x}-z)+\frac{1}{L}\|g_{x}\|^{2}+g_{x}^{\top}(z-x).
\end{align*}
Using that 
\[
f(z)\geq f(x)+\nabla f(x)^{\top}(z-x)+\frac{\mu}{2}\|z-x\|^{2}
\]
and that
\[
f(p_{x})\leq f(x)+\nabla f(x)^{\top}(p_{x}-x)+\frac{1}{2L}\|g_{x}\|^{2},
\]
we have the result:
\begin{align*}
\phi(z) & \geq h(p_{x})+f(x)+\nabla f(x)^{\top}(p_{x}-x)+\frac{1}{L}\|g_{x}\|^{2}+g_{x}^{\top}(z-x)+\frac{\mu}{2}\|z-x\|^{2}\\
 & \geq\phi(p_{x})+\frac{1}{2L}\|g_{x}\|^{2}+g_{x}^{\top}(z-x)+\frac{\mu}{2}\|z-x\|^{2}.
\end{align*}

\end{proof}
The next lemma shows that $\norm{g_{x}}_{2}\leq2G$ if $\phi$ is
$G$-Lipschitz.
\begin{lem}
\label{lem:Lipschitz_constant_g}If $\phi$ is $G$-Lipschitz, then
$\norm{g_{x}}_{2}\leq2G$ for all $x$.
\end{lem}

\begin{proof}
By the definition of gradient mapping (namely, $p_{x}$ is the minimizer
of a function), we have that
\[
f(x)+\nabla f(x)^{\top}(p_{x}-x)+\frac{L}{2}\norm{p_{x}-x}^{2}+h(p_{x})\leq f(x)+h(x).
\]
Using $h(p_{x})\geq h(x)+\nabla h(x)^{\top}(p_{x}-x)$, we have that
\[
0\geq\nabla\phi(x)^{\top}(p_{x}-x)+\frac{L}{2}\norm{p_{x}-x}^{2}\geq-G\norm{p_{x}-x}_{2}+\frac{L}{2}\norm{p_{x}-x}^{2}.
\]
Hence, we have that $\norm{p_{x}-x}_{2}\leq\frac{2}{L}G$ and hence
$\norm{g_{x}}_{2}\leq2G$.
\end{proof}

\subsection{Gradient Descent Using Gradient Mapping}

Putting $z=x$ in Theorem \ref{thm:gradient_mapping_lower_bound},
we have the following:
\begin{lem}[Gradient Descent Lemma]
\label{lem:gradient_descent_lemma}We have that
\[
\phi(p_{x})\leq\phi(x)-\frac{1}{2L}\norm{g_{x}}_{2}^{2}.
\]
\end{lem}

This shows that each step of the gradient step decreases the function
value by $\frac{1}{2L}\norm{g_{x}}_{2}^{2}$. Therefore, if the gradient
is large, then we decrease the function value by a lot. On the other
hand, Putting $z=x^{*}$ for Theorem \ref{thm:gradient_mapping_lower_bound}
shows that
\[
\phi(x^{*})\geq\phi(p_{x})+g_{x}^{\top}(x^{*}-x).
\]

If the gradient is small and domain is bounded, this shows that we
are close to the optimal. Combining these two facts, we can get the
gradient descent.

\subsection{Mirror Descent Using Gradient Mapping}

Consider the mirror descent as
\begin{equation}
x^{(k+1)}=x^{(k)}-\eta g_{x^{(k)}}.\label{eq:mirror_step}
\end{equation}
The main difference between this and gradient descent is that we will
take a larger step size $\eta$. To analyze the mirror descent, we
use Theorem \ref{thm:gradient_mapping_lower_bound} and the convexity
of $\phi$ to get that
\begin{equation}
\phi(\frac{1}{k}\sum_{i=1}^{k}p_{x^{(i)}})-\phi(x^{*})\leq\frac{1}{k}\sum_{i=1}^{k}(\phi(p_{x^{(i)}})-\phi(x^{*}))\leq\frac{1}{k}\sum_{i=1}^{k}g_{x^{(i)}}^{\top}(x^{(i)}-x^{*}).\label{eq:mirror_lemma}
\end{equation}
Therefore, it suffices to upper bound $g_{x^{(i)}}^{\top}(x^{(i)}-x^{*})$.
The following lemma shows that if $g_{x^{(i)}}^{\top}(x^{(i)}-x^{*})$
is large, then either the gradient is large or the distance to optimum
moves a lot. It turns out this holds for any vector $g$, not necessarily
an approximate gradient.
\begin{lem}[Mirror Descent Lemma]
\label{lem:mirror_step}Let $p=x-\eta g$. Then, we have that
\[
g^{\top}(x-u)=\frac{\eta}{2}\norm g_{2}^{2}+\frac{1}{2\eta}\left(\norm{x-u}_{2}^{2}-\norm{p-u}_{2}^{2}\right)
\]
for any $u$.
\end{lem}


\subsection{Algorithm and Analysis}

Recall Lemma \ref{lem:gradient_descent_lemma} shows that if the gradient
is large, gradient descent makes a large progress. On the other hand,
if the gradient is small, (\ref{eq:mirror_lemma}) shows that mirror
descent makes a large progress. Therefore, it is natural to combine
two approaches.

\begin{algorithm2e}[H]

\caption{$\mathtt{AGD}$}

\SetAlgoLined

\textbf{Input:} Initial point $x^{(1)}\in\Rn$.

$y^{(1)}\leftarrow x^{(1)}$, $z^{(1)}\leftarrow x^{(1)}$.

\For{$k=1,2,\cdots,T$}{

$x^{(k+1)}\leftarrow\tau z^{(k)}+(1-\tau)y^{(k)}$.

Perform a gradient step: $y^{(k+1)}=x^{(k+1)}-\frac{1}{L}g_{x^{(k+1)}}$.

Perform a mirror step: $z^{(k+1)}=z^{(k)}-\eta g_{x^{(k+1)}}$.

}

\textbf{Return $\frac{1}{T}\sum_{k=1}^{T}p_{x^{(k+1)}}$.}

\end{algorithm2e}

Note that if $\tau=1$, the algorithm is simply mirror descent and
if $\tau=0$, the algorithm is gradient descent.
\begin{thm}
\label{thm:AGD}Setting $\frac{1-\tau}{\tau}=\eta L$ and $\eta=\frac{1}{\sqrt{\mu L}}$,
we have that
\[
\phi(\frac{1}{T}\sum_{k=1}^{T}p_{x^{(k+1)}})-\phi(x^{*})\leq\frac{2}{T}\sqrt{\frac{L}{\mu}}\left(\phi(x^{(1)})-\phi(x^{*})\right).
\]
In particular, if we restart the algorithm every $4\sqrt{\frac{L}{\mu}}$
iterations, we can find $x$ such that 
\[
\phi(x)-\phi(x^{*})\leq2(1-\sqrt{\frac{\mu}{L}})^{\Omega(T)}\left(\phi(x^{(1)})-\phi(x^{*})\right)
\]
in $T$ steps. Furthermore, each step takes $\mathcal{T}_{f}+\mathcal{T}_{h,L}$
\end{thm}

\begin{proof}
Lemma \ref{lem:mirror_step} showed that
\begin{align}
g_{x^{(k+1)}}^{\top}(z^{(k)}-x^{*}) & \leq\frac{\eta}{2}\norm{g_{x^{(k+1)}}}_{2}^{2}+\frac{1}{2\eta}\left(\norm{z^{(k)}-x^{*}}_{2}^{2}-\norm{z^{(k+1)}-x^{*}}_{2}^{2}\right)\label{eq:AGD_1}
\end{align}

This shows that if the mirror descent has large error $g_{x^{(k+1)}}^{\top}(z^{(k)}-x^{*})$,
then the gradient descent makes a large progress ($\frac{\eta}{2}\norm{g_{x^{(k+1)}}}_{2}^{2}$).

To make the left-hand side usable, note that $x^{(k+1)}=z^{(k)}+\frac{1-\tau}{\tau}\cdot(y^{(k)}-x^{(k+1)})$
and hence
\begin{align*}
g_{x^{(k+1)}}^{\top}(x^{(k+1)}-x^{*}) & =g_{x^{(k+1)}}^{\top}(z^{(k)}-x^{*})+\frac{1-\tau}{\tau}\cdot g_{x^{(k+1)}}^{\top}(y^{(k)}-x^{(k+1)})\\
 & \leq g_{x^{(k+1)}}^{\top}(z^{(k)}-x^{*})+\frac{1-\tau}{\tau}(\phi(y^{(k)})-\phi(y^{(k+1)})-\frac{1}{2L}\norm{g_{x^{(k+1)}}}_{2}^{2})\\
 & \leq\frac{\eta}{2}\norm{g_{x^{(k+1)}}}_{2}^{2}+\frac{1}{2\eta}\left(\norm{z^{(k)}-x^{*}}_{2}^{2}-\norm{z^{(k+1)}-x^{*}}_{2}^{2}\right)+\frac{1-\tau}{\tau}(\phi(y^{(k)})-\phi(y^{(k+1)})-\frac{1}{2L}\norm{g_{x^{(k+1)}}}_{2}^{2}).
\end{align*}
where we used Theorem \ref{thm:gradient_mapping_lower_bound} in the
middle and (\ref{eq:AGD_1}) at the end.

Now, we set $\frac{1-\tau}{\tau}=\eta L$ and get
\[
g_{x^{(k+1)}}^{\top}(x^{(k+1)}-x^{*})\leq\eta L(\phi(y^{(k)})-\phi(y^{(k+1)}))+\frac{1}{2\eta}\left(\norm{z^{(k)}-x^{*}}_{2}^{2}-\norm{z^{(k+1)}-x^{*}}_{2}^{2}\right).
\]
Taking a sum on both side and let $\overline{x}=\frac{1}{T}\sum_{k=1}^{T}p_{x^{(k+1)}}$,
we have that
\begin{align*}
\phi(\overline{x})-\phi(x^{*}) & \leq\frac{1}{T}\sum_{k=1}^{T}(\phi(p_{x^{(k+1)}})-\phi(x^{*}))\\
 & \leq\frac{1}{T}\sum_{k=1}^{T}g_{x^{(k+1)}}^{\top}(x^{(k+1)}-x^{*})\\
 & \leq\frac{\eta L}{T}\left(\phi(y^{(1)})-\phi(y^{(T+1)})\right)+\frac{1}{2\eta T}\norm{z^{(1)}-x^{*}}_{2}^{2}\\
 & \leq\left(\frac{\eta L}{T}+\frac{1}{2\eta T\mu}\right)\left(\phi(x^{(1)})-\phi(x^{*})\right)
\end{align*}
The conclusion follows from our setting of $\eta$.
\end{proof}

\section{Accelerated Coordinate Descent}

Recall from Theorem \ref{thm:coordinate_descent} that if $\frac{\partial^{2}}{\partial x_{i}^{2}}\ell(x)\leq L_{i}$
for all $x$ and that $\ell$ is $\mu$ strongly convex, then we can
find $x^{(k)}$ in $k$ coordinate steps such that 
\[
\E\ell(x^{(k)})-\ell(x^{*})\leq(1-\frac{\mu}{L})^{\Omega(k)}(\ell(x^{(0)})-f(x^{*}))
\]
where $L=\sum_{i}L_{i}$. Now, the really fun part is here. Consider
the function 
\[
\phi(x)=f(x)+h(x)\quad\text{with}\quad f(x)=\frac{\mu}{2}\norm x_{2}^{2}\text{ and }h(x)=\ell(x)-\frac{\mu}{2}\norm x_{2}^{2}.
\]
Since $f$ is $\mu+\frac{L}{n}$ smooth (YES!, I know this is also
$\mu$ smooth) and $\mu$ strongly convex and since $h$ is convex,
we apply Theorem \ref{thm:AGD} and get an algorithm that takes $O^{*}(\sqrt{\frac{L}{n\mu}})$
steps. Note that each step involves $\mathcal{T}_{f}+\mathcal{T}_{h,\mu+\frac{L}{n}}$.
Obviously, $\mathcal{T}_{f}=0$. Next, note that $\mathcal{T}_{h,\mu+\frac{L}{n}}$
involves solving a problem of the form
\begin{align*}
y_{x} & =\argmin_{y}(\frac{\mu}{2}+\frac{L}{2n})\norm{y-x}^{2}+(\ell(y)-\frac{\mu}{2}\norm x^{2})\\
 & =\argmin_{y}\ell(y)-\mu y^{\top}x+\frac{L}{2n}\norm{y-x}^{2}.
\end{align*}
Now, we can apply Theorem \ref{thm:coordinate_descent} to solve this
problem. It takes 
\[
O^{*}(\frac{L+(L/n)\cdot n}{L/n})=O^{*}(n)\text{ coordinate steps.}
\]
Therefore, in total it takes
\[
O^{*}(\sqrt{\frac{L}{n\mu}})\cdot O^{*}(n)=O^{*}(\sqrt{\frac{Ln}{\mu}})\text{ coordinate steps}.
\]
Hence, we have the following theorem
\begin{thm}[Accelerated Coordinate Descent Convergence]
\label{thm:coordinate_descent_linear_acc}Given an $\mu$ strongly-convex
function $\ell$. Suppose that $\frac{\partial^{2}}{\partial x_{i}^{2}}\ell(x)\leq L_{i}$
for all $x$ and let $L=\sum L_{i}$. We can minimize $\ell$ in \textup{$O^{*}(\sqrt{\frac{Ln}{\mu}})$
coordinate steps.}
\end{thm}

\begin{rem}
It is known how to do it in $O^{*}(\sum_{i}\sqrt{\frac{L_{i}}{\mu}})$
steps \cite{allen2016even}.
\end{rem}


\section{Accelerated Stochastic Descent}

Recall from Theorem \ref{thm:stochastic_method} that if $\nabla^{2}\ell_{i}(x)\leq L$
for all $x$ and $i$ and that $\ell(x)=\frac{1}{n}\sum_{i=1}^{n}\ell_{i}(x)$
is $\mu$ strongly convex, then we can find $x$ in $O((n+\frac{L}{\mu})\log(\frac{1}{\epsilon}))$
stochastic steps such that 
\[
\E\ell(x)-\ell(x^{*})\leq\epsilon(\ell(x^{(0)})-\ell(x^{*})).
\]

Similar to the coordinate descent, we can accelerate it using the
accelerated gradient descent (Theorem \ref{thm:AGD}). To apply Theorem
\ref{thm:AGD}, we consider the function 
\[
\phi(x)=f(x)+h(x)\quad\text{with}\quad f(x)=\frac{\mu}{2}\norm x_{2}^{2}\text{ and }h(x)=\ell(x)-\frac{\mu}{2}\norm x_{2}^{2}.
\]
Since $f$ is $\mu+\frac{L}{n}$ smooth (YES!, I know this is also
$\mu$ smooth) and $\mu$ strongly convex and since $h$ is convex,
we apply Theorem \ref{thm:AGD} and get an algorithm that takes $O^{*}(1+\sqrt{\frac{L}{n\mu}})$
steps. Note that each step involves $\mathcal{T}_{f}+\mathcal{T}_{h,\mu+\frac{L}{n}}$.
Obviously, $\mathcal{T}_{f}=0$. Next, note that $\mathcal{T}_{h,\mu+\frac{L}{n}}$
involves solving a problem of the form
\begin{align*}
y_{x} & =\argmin_{y}(\frac{\mu}{2}+\frac{L}{2n})\norm{y-x}^{2}+(\ell(y)-\frac{\mu}{2}\norm x^{2})\\
 & =\argmin_{y}\ell(y)-\mu y^{\top}x+\frac{L}{2n}\norm{y-x}^{2}\\
 & =\argmin_{y}\frac{1}{n}\sum_{i}(\ell_{i}(y)+\frac{L}{2n}\norm{y-x}^{2}-\mu y^{\top}x)
\end{align*}
Now, we can apply Theorem \ref{thm:stochastic_method} to solve this
problem. It takes 
\[
O^{*}(n+\frac{L+\frac{L}{n}}{\mu+\frac{L}{n}})=O^{*}(n)
\]
Therefore, in total it takes
\[
O^{*}(1+\sqrt{\frac{L}{n\mu}})\cdot O^{*}(n)=O^{*}(n+\sqrt{n\frac{L}{\mu}}).
\]

\begin{thm}
\label{thm:acc_stochastic_method}Given a convex function $\ell=\frac{1}{n}\sum\ell_{i}$.
Suppose that $\nabla^{2}\ell_{i}(x)\leq L$ for all $i$ and $x$
and that $\ell$ is $\mu$ strongly convex. Suppose we can compute
$\nabla\ell_{i}$ in $O(1)$ time. We have an algorithm that outputs
an $x$ such that 
\[
\E\ell(x)-\ell(x^{*})\leq\varepsilon(\ell(x^{(0)})-\ell(x^{*}))
\]
in $O^{*}(n+\sqrt{n\frac{L}{\mu}})$ stochastic steps.
\end{thm}


